%% WMG.preliminary.tex
%% Windowed Minority Guidance: Preliminary Evidence for Timestep-Localized
%% Effects in Diffusion Denoising
%% Vittoria Lanzo — 2026
%% Single-scale (scale=1.0) EEML extended abstract — 4 pages

\documentclass[10pt]{article}

\usepackage{times}
\usepackage{amsmath}
\usepackage{amssymb}
\usepackage{booktabs}
\usepackage{hyperref}
\usepackage{microtype}
\usepackage{geometry}
\usepackage{listings}
\usepackage{xcolor}
\usepackage{multirow}
\usepackage{caption}

\geometry{
  letterpaper,
  top=1in, bottom=1in,
  left=1in, right=1in
}

\hypersetup{
  colorlinks=true,
  linkcolor=blue!70!black,
  citecolor=blue!70!black,
  urlcolor=blue!70!black
}

\lstset{
  basicstyle=\small\ttfamily,
  breaklines=true,
  frame=single,
  numbers=left,
  numberstyle=\tiny,
  xleftmargin=1.5em,
  framexleftmargin=1.5em
}

% -----------------------------------------------------------------------
\title{\textbf{Windowed Minority Guidance: Preliminary Evidence for
Timestep-Localized Effects in Diffusion Denoising}}

\author{%
  Vittoria Lanzo \\
  \textit{Independent Researcher}
}

\date{}
% -----------------------------------------------------------------------

\begin{document}

\maketitle

% -----------------------------------------------------------------------
\begin{abstract}
Minority guidance~\cite{um2024minority} steers diffusion sampling toward
under-represented data regions via a classifier gradient at every denoising
step.  We test whether this effect is timestep-localized, comparing three
equal-thirds windows---early $t\!\in\![0,333)$ (fine-detail, low noise),
mid $t\!\in\![333,667)$, late $t\!\in\![667,1000)$ (coarse structure, high
noise)---against full-chain and no-guidance baselines
over 50 seeds on LSUN Bedroom (250 runs).  Results confirm effects for mid,
early, and full conditions vs.\ baseline (all $p < 0.001$).  With late
vs.\ baseline narrowly missing correction ($p = 0.00716$ vs.\
$\alpha/7 = 0.00714$).  All windows produce detectable effects---preliminary
evidence for mid-chain dominance without strict localization.
\end{abstract}

% -----------------------------------------------------------------------
\section{Introduction}
\label{sec:intro}

Classifier-guided diffusion models~\cite{dhariwal2021diffusion} steer generation
with classifier gradients at every denoising step.  Um et al.\ \cite{um2024minority}
extend this to \emph{minority guidance}, directing sampling toward underrepresented
data regions.  Because decreasing timesteps produce qualitatively different effects
(Meng et al.\ \cite{meng2022sdedit}), the mid-chain phase may be most receptive to
minority gradient signals---enabling cheaper inference-time guidance if guidance
there recovers most of the full-chain effect.  We test this directly: Section~\ref{sec:discussion}
reports findings that partially challenge it.

% -----------------------------------------------------------------------
\section{Experiment}
\label{sec:experiment}

We use the LSUN Bedroom 256$\times$256 model with the frozen minority classifier of
Um et al.~\cite{um2024minority} (ImageNet feature extractor =
\texttt{mc\_lsun.pt}, \texttt{guidance\_scale=1.0},
\texttt{timestep\_respacing=250}).  Five conditions: \emph{baseline},
\emph{minority-full} $t\!\in\![0,1000)$, \emph{minority-early} $t\!\in\![0,333)$,
\emph{minority-mid} $t\!\in\![333,667)$, \emph{minority-late} $t\!\in\![667,1000)$.
Guidance is rescaled to the active window only; results may not
generalize.  We run 50 seeds $\times$ 5 conditions = 250 runs (Kaggle P100,
float32).  Seeds $1$--$10$ observed unbiased; results pre-registered in~\cite{um2024minority}.

\textbf{Primary metric}: classifier cross-entropy loss at class 99 (lower =
stronger guidance).  Softmax confidence is corroborated only---distributions are
severely right-skewed (mean/median: full 107x; late 154x; mid 40x; early 58x).
Switching from pre-registration classifier confidence to loss is an unplanned
procedural deviation; classifier-based effect sizes are one order of magnitude
smaller.

\paragraph{Timestep convention.}
We follow the DDPM convention~\cite{ho2020ddpm} in which $t=0$ corresponds to
the final denoising step (lowest noise, fine-detail refinement) and $t = T$ to
the first step (highest noise, coarse-structure formation).  The ``early''
window $t \in [0,333)$ therefore operates at \emph{low noise}; the ``late''
window $t \in [667, 1000)$ operates at \emph{high noise}.

\paragraph{Windowed \texttt{cond\_fn}.}
Guidance is applied exclusively within the designated timestep window.
Outside the window, the \texttt{cond\_fn} is a no-op, returning the unmodified
gradient zeroed out.  Full details are given in Appendix~\ref{sec:implementation}.

% -----------------------------------------------------------------------
\section{Results}
\label{sec:results}

\textbf{Table~\ref{tab:conditions}} ($n=50$; lower loss = stronger guidance).

\begin{table}[h]
\centering
\caption{Condition statistics ($n=50$ seeds; scale $= 1.0$; LSUN Bedroom).
Lower mean loss indicates stronger minority-class guidance.
Win Rate is the fraction of seeds on which the condition beats baseline.
Relative Effect $= (\text{baseline} - \overline{\text{cond}}) /
(\text{baseline} - \overline{\text{full}})$;
not a causal decomposition (windows interact across the rejection trajectory).}
\label{tab:conditions}
\renewcommand{\arraystretch}{1.1}
\begin{tabular}{@{}lcccc@{}}
\toprule
\textbf{Condition} & \textbf{Mean Loss} & \textbf{CV} &
  \textbf{Win Rate} & \textbf{Relative Effect} \\
\midrule
baseline        & 11.978 & 0.19 & ---   & 0.000 \\
minority-early  & 11.129 & 0.23 & 50/50 & 0.162 \\
minority-mid    &  9.591 & 0.35 & 39/50 & 0.456 \\
minority-late   & 10.717 & 0.30 & 34/50 & 0.241 \\
minority-full   &  6.741 & 0.61 & 44/50 & 1.000 \\
\bottomrule
\end{tabular}
\end{table}

The rank ordering mid $>$ late $>$ early $>$ baseline holds across 50 seeds
($r_\mathrm{rb} = 1.000, 0.812, 0.432, 0.925$ for early, mid, late, full
vs.\ baseline respectively).  \textbf{Table~\ref{tab:wilcoxon}} reports
Wilcoxon signed-rank tests, Bonferroni threshold $\alpha/7 = 0.0071$.

\begin{table}[h]
\centering
\caption{Wilcoxon signed-rank tests ($n=50$ pairs; \texttt{zero\_method=zsplit}).
$r_\mathrm{rb} \in [0,1]$: rank-biserial correlation (effect size);
$r_\mathrm{rb} > 0$ indicates first condition has lower loss.
Survives Bonferroni correction if $p < \alpha/7 = 0.0071$.}
\label{tab:wilcoxon}
\renewcommand{\arraystretch}{1.1}
\begin{tabular}{@{}llcc@{}}
\toprule
\textbf{Comparison} & \textbf{W} & \textbf{$p$-value} & \textbf{$r_\mathrm{rb}$} \\
\midrule
mid vs.\ baseline   & 120  & $<0.0001$ & 0.812 \\
late vs.\ baseline  & 362  & $0.0072$  & 0.432 \\
early vs.\ baseline &   0  & $<0.0001$ & 1.000 \\
full vs.\ baseline  &  48  & $<0.0001$ & 0.925 \\
mid vs.\ late       & 391  & $0.0166$  & 0.387 \\
late vs.\ early     & 547  & $0.3882$  & 0.142 \\
full vs.\ mid       & 181  & $<0.0001$ & 0.716 \\
\bottomrule
\end{tabular}
\end{table}

Mid recovers 0.456 of full-chain loss reduction (mean 9.591 vs.\ baseline
11.978 and full 6.741).  Pre-registered criterion: Bonferroni $p < 0.0071$.
Bootstrap 95\% CI on mid mean: [8.67, 10.50].  The relative effect
($\Delta_\mathrm{rel} = 0.456$) measures fraction of full-chain loss
reduction recovered; the CI bounds are absolute loss units and are not
directly comparable to $\Delta_\mathrm{rel}$.

\medskip

\noindent\textbf{Key finding}: early-window guidance (operating at low noise,
fine-detail phase) increases classifier confidence at every seed, demonstrating
that guidance applied to near-clean images reliably improves minority-class
quality ($r_\mathrm{rb} = 1.000$, $W = 0$, $p < 0.0001$).

\medskip

\noindent\textbf{Late $>$ Early at scale$=1.0$}: the mid-chain window separates
mid-chain from coarse-structure guidance; the coarse-structure (late) window
outperforms fine-detail (early) at this scale, though the difference is not
statistically significant ($p = 0.388$).

% -----------------------------------------------------------------------
\section{Discussion and Conclusion}
\label{sec:discussion}

Mid-window guidance recovers 0.456 of the full-chain effect---the largest
single-window contribution but insufficient alone.  All three windows produce
detectable seed-wise effects, contradicting binary localization.
\textbf{Key finding}: early-window guidance increases classifier-measured
guidance strength at every seed ($r_\mathrm{rb} = 1.000$, 50/50 wins)---evidence
that guidance applied to near-clean images reliably influences minority-class
sampling throughout the denoising trajectory.  Mid is quantitatively dominant
but the effect is distributed; mid alone recovers only 0.456 of the full-chain
loss reduction and cannot substitute for full-chain guidance at this scale.

\textbf{Limitations}: single configuration (LSUN Bedroom, class 99,
\texttt{guidance\_scale=1.0} vs.\ Um et al.'s scale $3.5$; $n=1$ sample per
seed); high full-chain variance (CV$=0.61$); no perceptual metrics (FID,
precision/recall).  Results may not generalize to other scales, datasets, or
architectures; the multi-scale sweep in the extended report examines scale
sensitivity directly.

\textbf{Revised framing}: mid is quantitatively dominant and the effect is
distributed; no single window approximates full-chain performance.  Late $>$
early is consistent with high-noise guidance producing conflicting gradients
before the denoising trajectory has committed to a minority-class layout.

% -----------------------------------------------------------------------
\bibliographystyle{plain}
\bibliography{references}

% -----------------------------------------------------------------------
\appendix
\section{Reproducibility}
\label{sec:reproducibility}

\subsection{Experimental Configuration}
\label{sec:config}

\paragraph{Models and checkpoints} (all publicly available):

\begin{table}[h]
\centering
\caption{Model checkpoints used in all experiments.}
\renewcommand{\arraystretch}{1.1}
\begin{tabular}{@{}lll@{}}
\toprule
\textbf{Component} & \textbf{Checkpoint} & \textbf{Source} \\
\midrule
Diffusion model     & \texttt{lsun\_bedroom.pt}     & OpenAI guided-diffusion \\
Feature extractor   & \texttt{256x256\_classifier.pt} & OpenAI guided-diffusion \\
Minority classifier & \texttt{mc\_lsun.pt}          & Um et al.\ (2024) \\
\bottomrule
\end{tabular}
\end{table}

\paragraph{Shared hyperparameters} (all fixed across conditions):

\begin{table}[h]
\centering
\caption{Hyperparameters shared across all conditions and seeds.}
\renewcommand{\arraystretch}{1.1}
\begin{tabular}{@{}ll@{}}
\toprule
\textbf{Parameter} & \textbf{Value} \\
\midrule
\texttt{diffusion\_steps}       & 1000 \\
\texttt{timestep\_respacing}    & 250 \\
\texttt{noise\_schedule}        & linear \\
\texttt{guidance\_scale}        & 1.0 \\
\texttt{image\_size}            & 256 \\
\texttt{num\_channels}          & 256 \\
\texttt{num\_res\_blocks}        & 2 \\
\texttt{attention\_resolutions} & 32,16,8 \\
\texttt{learn\_sigma}           & True \\
\texttt{use\_fp16}              & False \\
\texttt{manual\_class\_id}      & 99 \\
\texttt{num\_samples}           & 1 \\
\bottomrule
\end{tabular}
\end{table}

The minority classifier (\texttt{mc\_lsun.pt}) expects $8{\times}8$ latent features produced
by \texttt{256x256\_classifier.pt} with \texttt{pool=True}, \texttt{in\_channels=512},
\texttt{num\_channels=100}, \texttt{classifier\_width=128},
\texttt{classifier\_pool=attention}.

\textbf{unet.py patch}---the upstream repository's \texttt{UNetModel.forward}
is patched.  The following patch is required:

\begin{lstlisting}[language=Python,caption={unet.py patch: skip intermediate feature map during training.}]
def forward(self, x, timesteps, t_out=False):
    # skipping feature map extracted until final return
    if t_out:
        return h
    return self.out(h)  # original logic
\end{lstlisting}

\medskip
\noindent\textbf{Data schema} (\texttt{data/experiment\_runs.csv}; 250 rows, one per run):

\begin{center}
\small
\texttt{run\_id | condition | guidance\_window | t\_start | t\_end | guidance\_scale | seed | n\_generated | timestep | classifier\_mean\_loss | running\_seconds}
\end{center}

\noindent Metrics are computed at \texttt{t=0} (final denoising step) by passing
generated images through the frozen classifier chain.  Each run
produces one image; \texttt{classifier\_mean\_\*} fields are single-sample
values (\texttt{n=1} per seed per condition).

% -----------------------------------------------------------------------
\subsection{Windowed Guidance Implementation}
\label{sec:implementation}

The windowed guidance requires no architectural modification---it is
implemented entirely within the \texttt{cond\_fn} callback:

\begin{lstlisting}[language=Python,caption={Windowed \texttt{cond\_fn}: applies classifier gradient only within [t\_start, t\_end).}]
t_start = args.t_start   # e.g. 333 for mid window
t_end   = args.t_end     # e.g. 667 for mid window

def cond_func(x, t=None):
    current_t = t[0].item()           # actual denoising timestep
    # Zero grad outside the active window
    if not (t_start <= current_t < t_end):
        return torch.zeros_like(x)
    with torch.enable_grad():
        x_in = x.detach().requires_grad_(True)
        logits    = f_extractor(x_in, timesteps=t, t_out=True)
        log_probs = F.log_softmax(logits, dim=-1)
        selected  = log_probs[range(log_probs.shape[0]), y.view(-1)]
        # Scale applied only; sampler must NOT re-apply classifier_scale
        return torch.autograd.grad(selected.sum(), x_in)[0] \
               * args.classifier_scale
\end{lstlisting}

For baseline runs, \texttt{guidance\_scale=0} is passed rather than gating by
timestep.  Window boundaries are raw diffusion timesteps (0--1000); the active
window for mid is $t \in [333,667)$, corresponding to 334 of the 1000 raw
timesteps ($\approx\!134$ of the 250 respaced steps).

% -----------------------------------------------------------------------
\subsection{Statistical Analysis Code}
\label{sec:stats}

\begin{lstlisting}[language=Python,caption={Statistical analysis: Wilcoxon signed-rank tests and rank-biserial correlation (n=50).}]
import csv, random
from scipy import stats

# --- load paired losses ---
data = {}
with open('data/experiment_runs.csv') as f:
    for row in csv.DictReader(f):
        data.setdefault(row['condition'], {})[(int(row['seed']))] = \
            float(row['classifier_mean_loss'])

seeds = sorted(data['baseline'])
get   = lambda cond: [data[cond][s] for s in seeds]
baseline, early, mid, late, full = (get(c) for c in
    ['baseline','minority-early','minority-mid','minority-late','minority-full'])

# --- Wilcoxon signed-rank tests ---
def wilcoxon_rbc(a, b):
    diffs = [ai - bi for ai, bi in zip(a, b)]
    # zero_method='zsplit': split zero ranks equally between + and -
    W, p  = stats.wilcoxon(diffs, zero_method='zsplit')
    # rank-biserial: W / (n*(n+1)/2) * 2 - 1  (normalized to [0,1])
    n     = len(diffs)
    r_rb  = 1 - 2 * W / (n * (n + 1) / 2)
    return W, p, r_rb

comparisons = [
    (mid,   baseline, 'mid vs. baseline'),
    (late,  baseline, 'late vs. baseline'),
    (early, baseline, 'early vs. baseline'),
    (full,  baseline, 'full vs. baseline'),
    (mid,   late,     'mid vs. late'),
    (late,  early,    'late vs. early'),
    (full,  mid,      'full vs. mid'),
]

for a, b, label in comparisons:
    W, p, r_rb = wilcoxon_rbc(a, b)
    print(f'{label:<25}  W={W:>5.0f}  p={p:.4f}  r_rb={r_rb:.3f}')

alpha_corrected = 0.05 / 7  # = 0.00714

# --- Bootstrap 95% CI on mid mean (10,000 resamples) ---
random.seed(42)
boot  = sorted([sum(random.choices(mid, k=50)) / 50  for _ in range(10_000)])
ci    = [boot[249], boot[9749]]  # 2.5th and 97.5th percentiles

baseline_mean  = sum(baseline) / len(baseline)  # 11.978
threshold      = 0.05 / 7                        # 0.00714
mid_mean       = sum(mid) / len(mid)             # 9.591
#                                               (misses by 0.20)
\end{lstlisting}

\noindent\textbf{Environment}: Python 3.11, scipy 1.x, numpy 2.x.  All 250
runs generated on Kaggle (NVIDIA Tesla P100, float32, CUDA 11.x).  The Kaggle
notebook (including model download, patching, sampling loop, and CSV export)
is available on request.

\end{document}
